%!TEX root = ../main.tex
% APÉNDICE C — Detalles técnicos e implementación

\chapter{Detalles técnicos e implementación}\label{apx:tecnicos}

\begin{guia}
Todo lo que alguien necesita para reproducir los resultados:
hiperparámetros exactos, identificadores y versiones de los datos,
versiones de librerías, variables de entorno (sin secretos) y esquemas
de base de datos. Reproducibilidad total es un marcador de excelencia.
\end{guia}

\section{Versiones}

La tabla~\ref{tab:versiones} registra las versiones instaladas según el
archivo de bloqueo de dependencias del repositorio al 24 de septiembre de
2026, con las que se obtuvieron los resultados del
capítulo~\ref{cap:resultados}.

\begin{table}[H]
    \centering
    \caption{Versiones de las dependencias principales}
    \label{tab:versiones}
    \small
    \begin{tabular}{l l l}
        \toprule
        \tb{Componente} & \tb{Versión} & \tb{Uso} \\
        \midrule
        GraphSAST & 0.1.0 & Versión del motor, del CLI y del SARIF \\
        Node.js & 26.7.0 & Entorno de ejecución de la línea de comandos \\
        TypeScript & 5.9.3 & Lenguaje y compilador \\
        ts-morph & 23.0.0 & Análisis sintáctico \\
        Cytoscape.js & 3.34.0 & Visualización del grafo \\
        Vite & 5.4.21 & Empaquetado de la versión web \\
        Vitest & 2.1.9 & Pruebas automatizadas \\
        neo4j-driver & 6.2.0 & Cliente de Neo4j \\
        Neo4j & 5 Community & Base de datos de grafos (imagen \texttt{neo4j:5-community}) \\
        \bottomrule
    \end{tabular}
\end{table}

\section{Parámetros del análisis}

\begin{table}[H]
    \centering
    \caption{Parámetros del motor con sus valores por defecto}
    \label{tab:parametros}
    \small
    \begin{tabularx}{\textwidth}{l l L}
        \toprule
        \tb{Parámetro} & \tb{Valor} & \tb{Efecto} \\
        \midrule
        \texttt{-{}-max-depth} & 15 & Longitud máxima del camino entre origen y destino \\
        \texttt{-{}-max-file-bytes} & 1.000.000 & Tamaño máximo de un archivo analizado \\
        Análisis entre archivos & Activo & Se desactiva con \texttt{-{}-no-cross-file} \\
        \bottomrule
    \end{tabularx}
\end{table}

\section{Catálogo de detección}

El catálogo se compone de cuatro archivos JSON, uno por familia, en la
carpeta \texttt{catalog/} del repositorio. Cada archivo declara los
orígenes, los destinos sensibles y los saneamientos de la familia; su
contenido se resume en la tabla de debilidades del
capítulo~\ref{cap:marco-tecnologico}.

\section{Reproducción de los resultados}

\begin{lstlisting}[caption={Comandos para reproducir las mediciones del capítulo~\ref{cap:resultados}}, label={lst:reproduccion}]
npm install
npm test                 # suite automatizada (247 pruebas)
npm run benchmark        # banco sintetico (43 casos)
npm run cve:run          # advisories de npm
npm run scan -- scan docs/validacion/h3/caso-a --verbose
npm run baseline         # linea de base en el banco
npm run nodegoat         # NodeGoat, commit c5cb68a
\end{lstlisting}

El evaluador de NodeGoat clona el repositorio en el commit
\texttt{c5cb68a}, cuyo identificador completo está fijado en el
manifiesto \texttt{eval-nodegoat.json} junto con las vulnerabilidades
etiquetadas, el alcance del análisis y la regla de conteo. La copia se
guarda en el directorio temporal del sistema, o en el indicado por la
variable \texttt{NODEGOAT\_CACHE}.

\section{Variables de entorno}

El motor sobre Neo4j se habilita desde la API local
(\texttt{graphsast serve}) cuando está definida \texttt{NEO4J\_URI} y
apunta a la propia máquina; si apunta a otro equipo, se ignora y se usa
el recorrido en memoria, para que el código analizado no salga del
equipo. La instancia local se levanta con \texttt{npm run neo4j:up}.

\begin{table}[H]
    \centering
    \caption{Variables de entorno}
    \label{tab:variables-entorno}
    \small
    \begin{tabularx}{\textwidth}{l L l}
        \toprule
        \tb{Variable} & \tb{Propósito} & \tb{Valor por defecto} \\
        \midrule
        \texttt{NEO4J\_URI} & Dirección de la instancia de Neo4j; su presencia habilita el motor Cypher & Sin definir \\
        \texttt{NEO4J\_USER} & Usuario de la instancia & \texttt{neo4j} \\
        \texttt{NEO4J\_PASSWORD} & Contraseña de la instancia & La de desarrollo de \texttt{docker-compose.yml} \\
        \texttt{CVE\_JSON} & Archivo donde la validación con advisories escribe el informe en JSON & Sin definir \\
        \bottomrule
    \end{tabularx}
\end{table}
